\chapter*{ЗАКЛЮЧЕНИЕ}
\addcontentsline{toc}{chapter}{ЗАКЛЮЧЕНИЕ}

Таким образом, цель дипломной работы была достигнута --- разработан и реализован программно-алгоритмический
комплекс, снижающий вероятность несанкционированного копирования информации посредством управления доступом к \texttt{USB}-носителям.

Выполнены все поставленные задачи:
\begin{itemize}[label=---]
    \item проанализированы существующие решения в сфере управления доступом ко внешним запоминающим устройствам;
    \item формализованы требования к разрабатываемому комплексу с использованием диаграммы прецедентов;
    \item разработана архитектуру программно-алгоритмического комплекса;
    \item описаны компоненты и их взаимодействие;
    \item описаны ключевые структуры данных;
    \item обоснован выбор средств программной реализации;
    \item разработано программное обеспечение комплекса;
    \item проведено тестирования комплекса;
    \item описано взаимодействие пользователя с программным обеспечением;
    \item исследованы характеристики разработанного программного обеспечения.
\end{itemize}

Результаты исследования показали, что:

\begin{enumerate}
    \item количество актуальных записей о монтировании линейно влияет на время работы алгоритма актуализации записей;
    \item необходимость обновления информации в записи о монтировании увеличивает время инициализации приложения;
    \item необходимость перемонтирования устройств при инициализации приложения также увеличивает время инициализации.
\end{enumerate}